\section{Laufzeitanalyse}

Im Folgenden werden die diskrete Faltung in der Zeitdomäne und die Variante der Faltung,
basierend auf der Frequenzdomäne gegenübergestellt.

Dazu betrachten wir zunächst folgenden Pseudocode, der die diskrete Faltung in der
Zeitdomäne algorithmisch aufzeigt:

% Farben für Code definieren
\definecolor{codegreen}{rgb}{0,0.6,0}
\definecolor{codegray}{rgb}{0.5,0.5,0.5}
\definecolor{codepurple}{rgb}{0.58,0,0.82}
\definecolor{backcolour}{rgb}{0.95,0.95,0.92}

% Eigene Sprache für Pseudocode
\lstdefinelanguage{Pseudo}{
    morekeywords={
        algorithm, input, output, function, procedure,
        if, then, else, end, while, do, for, foreach,
        return, break, continue, repeat, until
    },
    sensitive=false,
    morecomment=[l]{//},
    morestring=[b]"
}

% Stil für Pseudocode
\lstdefinestyle{pseudocode}{
    language=Pseudo,
    backgroundcolor=\color{backcolour},
    commentstyle=\color{codegreen},
    keywordstyle=\color{blue},
    numberstyle=\tiny\color{codegray},
    stringstyle=\color{codepurple},
    basicstyle=\ttfamily\small,
    breakatwhitespace=false,
    breaklines=true,
    captionpos=b,
    keepspaces=true,
    numbers=left,
    numbersep=5pt,
    showspaces=false,
    showstringspaces=false,
    showtabs=false,
    tabsize=2,
    mathescape=true
}

\begin{lstlisting}[
    style=pseudocode,
    caption={Grundalgorithmus der Faltung im Zeitbereich},
    label={lst:directConvolution}
]
Algorithm directConvolution(inSignal, N, irSignal, M)
    outLen <- N + M - 1
    for i <- 0 to outLen
        for k <- 0 to M
            if i-k < 0 or n <= i-k
                continue
            outSignal[i] <- outSignal[i] + irSignal[k] * inSignal[i-k]
        done
    done
    return outSignal
\end{lstlisting}

Als Parameter erhält der Algorithmus der direkten Faltung zwei Listen.
Die erste Liste enthält die Folge der diskreten Abtastungssignale des Eingangssignals.
Dieses Signal hat den Betrag/ die Länge \textit{N}.
Die zweite Liste enthält die diskreten Abtastungssignale der Impulsantwort und hat die Länge \textit{M}.
Insgesamt hat die resultierende Folge eine Länge von $N + M - 1$.

Über diese Länge iteriert dementsprechend die äußere Schleife.
Die innere Schleife iteriert über die Länge $M$ der Impulsantwort.
Die if-Abfrage der inneren Schleife prüft, ob im Laufe der Iterationen
der Index des Eingangssignals 1 unterschreitet.
Durch diese Abzweigung an Sezenarien ergibt sich eine Best-Case Laufzeit von
\begin{equation}
\mathcal{O}((N + M - 1) \cdot (M + 2)) = \mathcal{O}(M^2 + NM + M + 2N - 2)
\end{equation}

Laut Polynomregel fallen Terme niederer Ordnung weg. Daraus folgt:
\begin{equation}
    \mathcal{O}(M^2 + MN) \implies \mathcal{O}(M \cdot N)
     , denn: N > M \implies N > M^2
\end{equation}

Die Laufzeitklasse der direkten, diskreten Convolution beträgt also laut O-Notation
\begin{equation}
\mathcal{O}(M \cdot N)
\end{equation}

